\section{Discussion and limitations}
The main conclusion is a classification of functional contact at original-mode rank-one points in capped TT classes. Cap-one cuts determine the exponent, while internal constraints can leave a singular feasible set without changing that exponent. If every cut permits an additional rank, all nonzero linear functionals return to first order. These statements describe the geometry of the declared class, not a universal benefit of tensor parameterization or rank selection.

The approximation result is consequential for interpretation. A confidence interval with an $n^{-3/2}$ pointwise length relies on a structural allowance of the same order if that gain is to persist uniformly. An ambient coefficient departure of order between $n^{-3/2}$ and $n^{-1/2}$ can be statistically undetectable and still dominate the target uncertainty. The theorem therefore argues for explicit sensitivity analysis whenever rank restrictions are approximate. It does not provide a way to infer the required small allowance from the same sample without additional assumptions.

The limitations are substantial. The pointwise length theorem fixes ambient dimension and the signal directions, with an original-mode rank-one evaluation point. The OLS ball can incur large dimension constants, and the experiments show that the computable interval is far from a sharp constant procedure. No matching dimension-dependent minimax length, efficient confidence constant, or general classification at non-rank-one TT points has been established. The inherited hardness of an exact local coefficient does not establish hardness of statistically optimal inference.

\paragraph{Status of possible extensions.}
\textsc{Partial}: growing-dimensional inference has a valid ambient localization construction and fixed-geometry lower bounds, but their dimension factors do not match. \textsc{Conjectured}: a general classification at non-rank-one points with several weak bonds may require higher-order feasible arcs beyond the present elementary paths; no universal rate is asserted. \textsc{Partial}: the spectral-norm reduction identifies a difficult constant but leaves computationally optimal intervals unresolved. \textsc{Conjectured}: non-Gaussian or dependent-noise versions would need new honest localization and moment bounds. These statements are research questions, not theorems of this paper.

\subsection{Relevance to modern AI and why no LLM benchmark is claimed}
A tensor regression head on frozen features offers a direct statistical setting for coefficient functionals. In a locally linear adapter model, a fixed evaluation contrast has the form $\ip{A}{B}$, with $A$ determined by held-out or externally fixed Jacobian features. Existing TT adapters make such restrictions plausible modeling choices, but a trained LLM does not automatically satisfy isotropic design, Gaussian residuals, fixed dimensions, or a known approximation allowance. A fully observed checkpoint is also a deterministic approximation object, not the sampling experiment in \eqref{eq:model}.

Accordingly, this draft contains no real-data or trained-LLM result. The companion application protocol specifies a frozen-feature regression study with a separately fixed functional, sample splitting for any learned loading, residual checks, and an explicit sensitivity grid. A convincing application would need to justify the population target and structural approximation bound, rather than add an LLM benchmark to support a general mathematical claim. The present theory applies independently of any AI interpretation.

\paragraph{Research assessment.}
The completed fixed-dimensional results support a focused theoretical manuscript for further expert evaluation. The central novelty candidate is the cap-profile classification and its inferential interpretation, not generic higher-order asymptotics. A broader theory of local moduli at singular structured parameters, or dimension-dependent optimality with a useful procedure, would be needed to make a strong case for the Annals of Statistics. The current draft should not be treated as submission ready solely because its proofs and experiments are complete.

\paragraph{Reproducibility and proof status.}
The accompanying project includes the full literature matrix, theorem ledger, assumption audit, canonical proof source, fixed-seed simulation code, replicate outputs, vector figures, and a workbook containing the plotted summaries and source metadata. Results marked \pr\ have complete derivations and internal mathematical checks. This label does not certify priority or replace independent mathematical review. The empirical addendum and the later supporting local-Wald derivation are disclosed in the project records.
